\documentclass[11pt,a4paper]{article}

\usepackage[margin=2.5cm]{geometry}
\usepackage{amsmath,amssymb,mathtools,amsthm}
\usepackage{booktabs,array}
\usepackage{algorithm}
\usepackage{algpseudocode}
\usepackage[dvipsnames]{xcolor}
\usepackage{microtype}
\usepackage[colorlinks=true,linkcolor=MidnightBlue,citecolor=MidnightBlue,urlcolor=MidnightBlue]{hyperref}
\usepackage{cleveref}

\newtheorem{proposition}{Proposition}
\theoremstyle{remark}\newtheorem{remark}{Remark}
\newcommand{\R}{\mathbb{R}}
\newcommand{\dd}{\,\mathrm{d}}
\newcommand{\cR}{c_{\mathrm{ref}}}
\newcommand{\SolveM}{\mathcal{M}}
\DeclareMathOperator{\diag}{diag}

\title{The nonlinear Poisson--Boltzmann equation on a twisted elliptic torus:\\
two fixed-operator iterations for an OMS / HBS solver}
\author{}
\date{}

\begin{document}
\maketitle

\begin{abstract}
We state the normalized nonlinear Poisson--Boltzmann (PB) equation for a pure
electron plasma confined on magnetic surfaces, pose it on a twisted elliptic
torus with prescribed profiles, and construct a manufactured solution for
verification. We then describe two iterations that solve it while keeping a
\emph{single} linear operator $M=-\Delta+\cR(x)$: a fixed-point
iteration on the OMS interface system, whose state is only the interface
traces, and a volumetric Newton--Krylov method preconditioned by $M$. In both, the
slab factorizations, the hierarchical (HBS) compressions and the red--black
preconditioner of the overlapping multi-slab (OMS) solver are built once; every
iteration only changes right-hand sides and pointwise (diagonal) volume
coefficients. No slab factorization is repeated.
\end{abstract}

\tableofcontents

%======================================================================
\section{The model}
%======================================================================

\subsection{Physical equation}
For a pure electron plasma confined on magnetic surfaces (Costello \& Helander,
following Pedersen \& Boozer), the electrostatic potential $\Phi$ obeys
Poisson's equation with the electron charge density,
\begin{equation}
  \varepsilon_0\nabla^2\Phi = e\,n_e ,
\end{equation}
and force balance along the magnetic field, with isothermal flux surfaces
$T_e=T_e(\psi)$, gives a Boltzmann relation on each surface,
\begin{equation}
  n_e = N(\psi)\exp\!\Big(\frac{e\Phi}{T_e(\psi)}\Big).
\end{equation}
Here $e>0$ is the elementary charge, $\psi$ labels the flux surfaces and
$N(\psi)$ is a flux function. Combining,
\begin{equation}\label{eq:pb-dim}
  \nabla^2\Phi = \frac{e}{\varepsilon_0}\,N(\psi)\exp\!\Big(\frac{e\Phi}{T_e(\psi)}\Big).
\end{equation}
(The paper prints this as $\nabla^2\Phi=N(\psi)\exp(e\Phi/T_e)$; the factor
$e/\varepsilon_0$ follows from the Poisson equation stated with it.)

\subsection{Normalization}
With a length scale $a$, a reference temperature $T_{\mathrm{ref}}$ and a
reference density $n_{\mathrm{ref}}$, set
\begin{equation}
  x = a\hat x,\qquad \phi=\frac{e\Phi}{T_{\mathrm{ref}}},\qquad
  \hat N=\frac{N}{n_{\mathrm{ref}}},\qquad \tau=\frac{T_e}{T_{\mathrm{ref}}},\qquad
  \lambda_D=\sqrt{\frac{\varepsilon_0T_{\mathrm{ref}}}{e^2n_{\mathrm{ref}}}},\qquad
  \Lambda=\frac{a}{\lambda_D}.
\end{equation}
Then \eqref{eq:pb-dim} becomes (dropping hats)
\begin{equation}\label{eq:pb}
  \boxed{\;-\Delta\phi + \Lambda^2 N(\psi)\,e^{\phi/\tau(\psi)} = 0\quad\text{in }\Omega,
  \qquad \phi=\phi_b\ \text{on }\partial\Omega.\;}
\end{equation}
$\Lambda$ is the number of Debye lengths across the device; the regime of
interest is $\Lambda\gg1$, where the exponential term is stiff.

We use \emph{model (A)}: $\Lambda$, the flux label $\psi(x)$, and the profiles
$N(\psi)>0$, $\tau(\psi)>0$ are prescribed, with Dirichlet data on the plasma
boundary. (The alternative, prescribing the surface-averaged density
$\langle n_e\rangle_\psi$ and solving for $N$, makes the equation nonlocal on
each surface; it is not considered here.)

\paragraph{Notation for the nonlinearity.} The nonlinear term is a pointwise
function of the position $x\in\Omega$ and of the \emph{scalar} value $s\in\R$
of the potential there:
\begin{equation}\label{eq:g}
  g(x,s)=\Lambda^2N(\psi(x))\,e^{s/\tau(\psi(x))},\qquad
  \partial_sg(x,s)=\frac{g(x,s)}{\tau(\psi(x))}>0 .
\end{equation}
For a field $\phi$ on $\Omega$ we write the nonlinear term and the coefficient
of its linearization as fields,
\begin{equation}\label{eq:Gc}
  G[\phi](x)=g\big(x,\phi(x)\big),\qquad c[\phi](x)=\partial_sg\big(x,\phi(x)\big),
  \qquad G[\phi+\delta]=G[\phi]+c[\phi]\,\delta+\mathcal O(\delta^2).
\end{equation}
Only $G[\cdot]$ and $c[\cdot]$ appear below, except in contraction factors,
where $\partial_sg$ is evaluated at intermediate values. With this notation
\eqref{eq:pb} reads $-\Delta\phi+G[\phi]=0$. Since $g$ is increasing in $s$, it
is the Euler--Lagrange equation of a strictly convex energy and has exactly one
solution.

\subsection{Maximum-principle estimate}
The following elementary bound is used repeatedly.
\begin{proposition}\label{prop:maxp}
Let $c\ge0$ and let $w$ solve $-\Delta w+c(x)w=h$ in $\Omega$, $w=0$ on $\partial\Omega$. Then
\[
\|w\|_\infty\le C_P\,\|h\|_\infty,\qquad C_P=\max_\Omega w_0,\ \ -\Delta w_0=1,\ w_0|_{\partial\Omega}=0,
\]
and, if $c>0$, also $\|w\|_\infty\le\|h/c\|_\infty$. If $\Omega$ lies between
two parallel planes a distance $d$ apart, $C_P\le d^2/8$.
\end{proposition}
\begin{proof}[Sketch]
At a positive interior maximum $x_0$ of $w$, $-\Delta w(x_0)\ge0$, so
$c\,w(x_0)\le h(x_0)$; the same at a negative minimum. The first bound follows
by comparison with $\|h\|_\infty w_0$, and $w_0$ is dominated by the
one-dimensional torsion function of the slab.
\end{proof}

%======================================================================
\section{Geometry}
%======================================================================
The domain is a twisted elliptic torus, parametrized over a straight periodic
cylinder $C=[0,L)\times D_{R_c}$, $D_{R_c}$ the disk of radius $R_c$. With
reference coordinates $p=(t,x,y)$, major radius $R_0$, ellipse semi-axes
$a_e,b_e$ and an integer twist $\mathrm{tw}$,
\begin{equation}
\begin{gathered}
  R_c=\sqrt{a_eb_e},\quad L=2\pi R_0,\quad \zeta=\frac{t}{R_0},\quad
  X=\frac{a_e}{R_c}x,\quad Y=\frac{b_e}{R_c}y,\quad \alpha=\mathrm{tw}\,\zeta,\\
  R=R_0+X\cos\alpha-Y\sin\alpha,\qquad
  \chi(p)=\big(R\cos\zeta,\;R\sin\zeta,\;X\sin\alpha+Y\cos\alpha\big).
\end{gathered}
\end{equation}
The map preserves the centreline length and the cross-sectional area. Its
inverse $y=\chi^{-1}$ and the derivatives $\nabla y_k$, $\Delta y_k$ are
available in closed form. In reference coordinates the Laplacian becomes
\begin{equation}
  -\Delta u = -\sum_{k,l}\nabla y_k\!\cdot\!\nabla y_l\;\partial_k\partial_l\tilde u
  \;-\;\sum_k(\Delta y_k)\,\partial_k\tilde u,\qquad \tilde u=u\circ\chi,
\end{equation}
and zero-order terms and body loads are simply evaluated at $\chi(p)$. We use
the toy flux label
\begin{equation}
  \psi(\chi(t,x,y))=\frac{x^2+y^2}{R_c^2}\in[0,1],
\end{equation}
whose level sets are nested twisted elliptic tori, $\psi=1$ being the plasma
boundary.

The cylinder is discretized by a five-patch high-order (HPS) discretization of
the cross-section (a square core and four arched patches), extruded along $t$,
and the solver is the overlapping multi-slab method (OMS): $t$ is split into
$N_s$ double slabs, the unknowns are the traces on the slab midplanes, and the
reduced system $(I+S)u=r$ couples neighbouring interfaces through the local
solution maps $S$. The blocks of $S$ are compressed as HBS matrices, and a
red--black (cyclic) reduction of the compressed system serves as preconditioner.

%======================================================================
\section{Manufactured solution}
%======================================================================
Choose a smooth $\phi^*$ in physical coordinates and define the source
\begin{equation}
  f^*= -\Delta\phi^*+G[\phi^*].
\end{equation}
Then $\phi^*$ is the unique solution of
\begin{equation}\label{eq:mms}
  F(\phi):=-\Delta\phi+G[\phi]-f^*=0\ \text{in }\Omega,\qquad \phi=\phi^*\ \text{on }\partial\Omega .
\end{equation}
A convenient choice has a closed-form Laplacian:
\begin{equation}
  \phi^*(x)=A\,\sin(\kappa_1x_1+\theta)\cos(\kappa_2x_2)\,e^{\kappa_3x_3}+B,
  \qquad -\Delta\phi^*=(\kappa_1^2+\kappa_2^2-\kappa_3^2)\,(\phi^*-B),
\end{equation}
with profiles, e.g.,
\begin{equation}
  N(\psi)=1-\tfrac12\psi,\qquad \tau(\psi)=1-0.4\,\psi .
\end{equation}
Pick $A$ so that $\phi^*/\tau$ varies by a few units over $\Omega$; then
the coefficient $c[\phi^*]$ varies by a factor $e^{\mathcal O(\text{few})}$ and the nonlinearity matters.
Sweep $\Lambda\in\{1,10,30\}$. The body load is $f^*\circ\chi$ on reference
points, and the boundary data $\phi^*\circ\chi$ on the tube wall $x^2+y^2=R_c^2$.

%======================================================================
\section{The fixed reference operator}
%======================================================================
\subsection{What is built once}
Fix a smooth $\cR(x)\ge0$ and the linear operator
\begin{equation}
  M=-\Delta+\cR(x).
\end{equation}
For $M$ the OMS solver builds, once:
\begin{enumerate}
  \item per double slab, the sparse factorization (MUMPS) of the reduced local system $A_{ii}$;
  \item the interface maps $S$ (applied exactly through these factorizations) and their HBS compressions;
  \item the red--black preconditioner from the compressed blocks.
\end{enumerate}
In terms of these, define the volumetric solve
\begin{equation}
  w=\SolveM(h,b)\quad\Longleftrightarrow\quad Mw=h\ \text{in }\Omega,\ \ w=b\ \text{on }\partial\Omega ,
\end{equation}
carried out as: reduce the body load $h$ on every slab (leaf-level static
condensation), form the local right-hand sides with the kept factorizations,
solve the interface system by GMRES preconditioned by red--black, and
reconstruct $w$ on the leaf grids (kept factorizations plus leaf solves);
the slab-level steps are those of \Cref{prop:condense}. $\SolveM$ involves no factorization.

\subsection{Why the coefficient cannot be changed on the interface system}
The interface blocks are $S=-(A_{ii}^{-1}A_{ib})|_{\text{midplane},\,\text{faces}}$.
The coefficient $c$ enters every leaf map and $A_{ii}^{-1}$, so $S$ depends on
$c$ nonlinearly and non-locally: changing $c$ is \emph{not} a diagonal shift of
the interface operator. In one dimension, for $-u''+cu=0$ on a double slab of
half-width $H$,
\begin{equation}
  u(0)=\frac{u(-H)+u(H)}{2\cosh(\sqrt c\,H)}\qquad\text{versus}\qquad \frac{u(-H)+u(H)}{2}\ \ (c=0).
\end{equation}
A new $c$ would therefore require refactoring and recompressing every slab. On
the \emph{volumetric} operator, however, $-\Delta+c$ is the Laplacian plus a
diagonal. Both approaches below exploit that while keeping $S$ fixed: the
interface iteration through local slab solves with the kept factorizations,
Newton--Krylov through $\SolveM$.

%======================================================================
\section{Approach 1: fixed-point iteration on the interface system}\label{sec:iface}
%======================================================================
The state of this iteration is only the vector of interface traces
$u=(u_0,\dots,u_{N_s-1})$ on the slab midplanes $\Gamma_i$, the same unknowns
as the linear OMS system. Volumetric fields exist only transiently, inside
one slab at a time.

\subsection{Nonlinear interface system}
Let $\Omega_i$ be double slab $i$. Its end faces are the neighbouring midplanes
$\Gamma_{i-1}$ and $\Gamma_{i+1}$, its lateral boundary $W_i$ is part of the
tube wall, and its midplane is $\Gamma_i$. Given the neighbours' traces, the
\emph{local nonlinear problem} is
\begin{equation}\label{eq:local}
  -\Delta\phi^{(i)}+G[\phi^{(i)}]=f^*\ \text{in }\Omega_i,\qquad
  \phi^{(i)}=u_{i\mp1}\ \text{on }\Gamma_{i\mp1},\qquad \phi^{(i)}=\phi^*\ \text{on }W_i .
\end{equation}
By monotonicity it has exactly one solution. Define the local trace maps $\mathcal E_i$ (solve on $\Omega_i$, then restrict
to $\Gamma_i$) and the interface residual
\begin{equation}
  \mathcal E_i(u)=\phi^{(i)}\big|_{\Gamma_i},\qquad
  \mathcal R(u)=u-\mathcal E(u).
\end{equation}

\begin{proposition}\label{prop:equiv}
$\mathcal R(u)=0$ if and only if $u$ is the trace on the midplanes of the
solution $\phi^*$ of \eqref{eq:mms}.
\end{proposition}
\begin{proof}[Sketch]
The traces of $\phi^*$ satisfy \eqref{eq:local} on every slab, so $\mathcal R=0$.
Conversely, let $\mathcal R(u)=0$ and consider the overlap cell between
$\Gamma_i$ and $\Gamma_{i+1}$, which belongs to both $\Omega_i$ and
$\Omega_{i+1}$. On it, $\phi^{(i)}$ and $\phi^{(i+1)}$ solve the same equation
with the same wall data and the same values $u_i$ and $u_{i+1}$ at its two
ends. By uniqueness they coincide, so the local solutions patch together into
a global solution, which is $\phi^*$.
\end{proof}

This is the nonlinear counterpart of the linear OMS system, and in the
discrete setting the same argument applies slab by slab. Linearizing the local problems shows that
\begin{equation}
  \mathcal E'(u)=-S\big(c[\phi_u]\big),
\end{equation}
where, for a coefficient field $c\ge0$, $S(c)$ denotes the OMS interface
operator of the \emph{linear} operator $-\Delta+c$, and $\phi_u$ is the field
assembled from the local solutions for the data $u$. So the
linear OMS system built once for $c=\cR$, $I+S_{\mathrm{ref}}=I+S(\cR)$, is a
frozen Jacobian of $\mathcal R$.

\subsection{Iteration}
With $P\approx(I+S_{\mathrm{ref}})^{-1}$, for example one application of the
red--black preconditioner or $(I+S_{\mathrm{ref}})^{-1}$ applied by
preconditioned GMRES,
\begin{equation}\label{eq:ichord}
  \boxed{\;u_{k+1}=u_k-P\,\mathcal R(u_k).\;}
\end{equation}
The fixed point is exact for any choice of $P$; $P$ only sets the rate.
Linearizing about the solution,
\begin{equation}\label{eq:irate}
  e_{k+1}\approx\big[I-P\,(I+S(c_*))\big]e_k
  =P\Big[\big(P^{-1}-(I+S_{\mathrm{ref}})\big)+\big(S_{\mathrm{ref}}-S(c_*)\big)\Big]e_k,
  \qquad c_*=c[\phi^*].
\end{equation}
The rate is therefore governed by two quantities:
\begin{itemize}
  \item the accuracy of the preconditioner, $P^{-1}\approx I+S_{\mathrm{ref}}$;
  \item the mismatch of the linearized local maps, $S(\cR)$ against $S(c_*)$.
\end{itemize}
The interface maps measure how boundary data propagates to the midplane.
Under strong screening ($\Lambda\gg1$, $\lambda_D\ll H$) both decay like
$e^{-\sqrt{c_{\min}}\,H}$, so $I+S$ is close to the identity, the outer
iteration needs few steps, and almost all the work is local.

\subsection{Local nonlinear solves in statically condensed form}\label{sec:local}
$\mathcal E_i(u_k)$ is evaluated on slab $i$ alone. What is kept for the slab
is not the volumetric operator $M_i=-\Delta+\cR$ on $\Omega_i$ but its static
condensation onto the leaf skeleton, together with the leaf operators that
produced it. We first write the application of $M_i^{-1}$ in these terms, then
the local iteration.

\paragraph{Leaf operators.} Let $\mathcal L_i$ be the leaves of slab $i$. On
leaf $\ell$ the unknowns are the values $\phi_{I,\ell}$ at interior collocation
nodes and $x_\ell$ at face nodes. Write the leaf discretization of $-\Delta$ as
$[K_{II,\ell}\ K_{IF,\ell}]$ (rows at interior nodes) and let $C_\ell$ be the
diagonal of $\cR$ at the interior nodes, so the leaf rows of $M_i$ are
$[K_{II,\ell}+C_\ell\ \ K_{IF,\ell}]=:[L_{II,\ell}\ L_{IF,\ell}]$. Let $D_\ell$
map a leaf field $(\phi_{I,\ell},x_\ell)$ to its outgoing (conormal) flux on
the leaf faces. Define
\begin{equation}\label{eq:leafops}
  Y_\ell=-L_{II,\ell}^{-1}L_{IF,\ell},\qquad
  T_\ell=D_\ell\begin{bmatrix}Y_\ell\\ I\end{bmatrix},\qquad
  q_\ell(h)=D_\ell\begin{bmatrix}L_{II,\ell}^{-1}h_\ell\\ 0\end{bmatrix}:
\end{equation}
$Y_\ell$ extends face data into the leaf, $T_\ell$ is the leaf
Dirichlet-to-Neumann map, and $q_\ell(h)$ is the flux of the particular
solution for the body load $h_\ell$ (the restriction of $h$ to the interior
nodes of $\ell$) with zero face data.

\paragraph{Skeleton system.} With $E_\ell$ the restriction from the slab
skeleton to the faces of $\ell$, flux continuity on every skeleton face gives
\begin{equation}\label{eq:skel}
  A\,x=-Q\,h,\qquad A=\sum_\ell E_\ell^{\mathsf T}T_\ell E_\ell,\qquad
  Q\,h=\sum_\ell E_\ell^{\mathsf T}q_\ell(h).
\end{equation}
Splitting the skeleton into interior faces ($I$) and faces on
$\partial\Omega_i=\Gamma_{i-1}\cup\Gamma_{i+1}\cup W_i$ ($b$), where the
Dirichlet data $b_i=(u_{i-1},u_{i+1},\phi^*|_{W_i})$ is prescribed,
$A_{ii}x_I=-(Qh)_I-A_{ib}b_i$. $A_{ii}$ is the matrix whose factorization is
kept. (For the five-patch cylinder, the leaf flux is converted to the physical
conormal flux and the patches are glued before this split; the body-load
reduction must use the same conversion, or it applies a different $M_i$.)

\begin{proposition}\label{prop:condense}
For any body load $h$, the three steps
\begin{enumerate}
  \item \emph{reduce:} $\beta=-(Qh)_I-A_{ib}b_i$ (one batch of leaf interior solves);
  \item \emph{skeleton solve:} $x_I=A_{ii}^{-1}\beta$ (kept factorization), $x_b=b_i$;
  \item \emph{reconstruct:} $\phi|_\ell=Y_\ell E_\ell x+L_{II,\ell}^{-1}h_\ell$ (one batch of leaf interior solves)
\end{enumerate}
return exactly the solution $\phi$ of $M_i\phi=h$ with Dirichlet data $b_i$.
\end{proposition}
\begin{proof}[Sketch]
Static condensation is block Gaussian elimination of the interior nodes of each
leaf, followed by elimination of the interior skeleton; both are exact.
\end{proof}

\paragraph{Local chord iteration.} Hence the chord iteration
$\phi_{j+1}=M_i^{-1}(f^*-G[\phi_j]+\cR\phi_j)$ on $\Omega_i$ (with the data of
\eqref{eq:local}) is carried out as
\begin{equation}\label{eq:lchord}
\boxed{\;
\begin{aligned}
  h_j&=f^*-G[\phi_j]+\cR\,\phi_j &&\text{pointwise at the leaf interior nodes},\\
  x_{j+1,I}&=A_{ii}^{-1}\big(-(Qh_j)_I-A_{ib}b_i\big),\quad x_{j+1,b}=b_i,\\
  \phi_{j+1}|_\ell&=Y_\ell E_\ell\,x_{j+1}+L_{II,\ell}^{-1}h_{j,\ell}
  &&\text{for every }\ell\in\mathcal L_i .
\end{aligned}\;}
\end{equation}
(Superscripts $(i)$ are dropped.) By \Cref{prop:condense} this is the same
iteration as applying $M_i^{-1}$, so its fixed point and its rate are
unchanged. The error $e_j=\phi_j-\phi^{(i)}$ satisfies
$(M_ie_{j+1})(x)=\big(\cR(x)-\partial_sg(x,\xi_j(x))\big)\,e_j(x)$ with zero
Dirichlet data, for some $\xi_j(x)$ between $\phi_j(x)$ and $\phi^{(i)}(x)$
(mean value theorem). If $\cR>0$, \Cref{prop:maxp} on $\Omega_i$ gives
contraction with factor
\begin{equation}\label{eq:rho}
  \rho_i=\sup_{x\in\Omega_i,\ s\in I(x)}\Big|1-\frac{\partial_sg(x,s)}{\cR(x)}\Big|,
\end{equation}
where $I(x)$ is the interval spanned by the iterates and the solution at $x$.
(The supremum must be restricted to such an interval: over all $s\in\R$ it is
infinite.) \Cref{prop:maxp} is a statement about the continuous problem; for
the spectral discretization it holds up to discretization error for resolved
fields, and the observed rate should be checked against \eqref{eq:rho}.

\paragraph{What one local step costs and needs.}
\begin{itemize}
  \item Two batches of small dense leaf solves (reduction and reconstruction)
    and one solve with the kept MUMPS factorization. No factorization of any
    kind.
  \item The reconstruction cannot be skipped: $G[\phi_j]$ is needed at every
    leaf interior node, so the state of the local iteration is the
    collocation field on the slab's leaves, not the skeleton values alone.
    The midplane trace $\mathcal E_i(u)$, by contrast, is read directly from
    the skeleton solution $x$ on the faces of $\Gamma_i$.
  \item Per leaf, the factors of $L_{II,\ell}$ (or $L_{II,\ell}^{-1}$), $Y_\ell$ and
    $D_\ell$ must be kept, built with the same $\cR$ as $A_{ii}$. They are small,
    but in the present slab builder they are discarded after gluing.
  \item The slabs are independent, so they can be processed in parallel or one
    at a time, each with one transient volumetric field. The previous outer
    iteration's local field, if kept, is a good starting guess. The local
    tolerance can be loose in early outer iterations and tightened as
    $\mathcal R(u_k)\to0$ (inexact evaluation of $\mathcal R$).
\end{itemize}

\begin{remark}[Nonlinear elimination at leaf level]\label{rem:leafelim}
Because leaf solves are small and dense, the leaves can absorb the
nonlinearity exactly, leaving a frozen condensed operator only on the
skeleton. For face data $x$, let $\phi_{I,\ell}(x)$ solve the nonlinear leaf
problem
\begin{equation}
  K_{II,\ell}\,\phi_{I,\ell}+K_{IF,\ell}\,E_\ell x+G[\phi]_{I,\ell}=f^*_{I,\ell},
\end{equation}
by a few Newton steps with the exact leaf Jacobian $K_{II,\ell}+\diag(c[\phi]_{I,\ell})$
(new small dense factorizations, batched, no MUMPS), warm-started from the
previous leaf field. This defines nonlinear leaf fluxes
$\mathcal T_\ell(x)=D_\ell\big(\phi_{I,\ell}(x),E_\ell x\big)$ and the skeleton
residual
\begin{equation}
  \mathcal R_{\mathrm{skel}}(x_I)=\Big(\sum_\ell E_\ell^{\mathsf T}\mathcal T_\ell(x)\Big)_I,
  \qquad x_b=b_i ,
\end{equation}
whose Jacobian is the condensed operator of the linearization,
$A_{ii}(c[\phi_x])$. The kept $A_{ii}=A_{ii}(\cR)$ is a frozen Jacobian, and
\begin{equation}
  x_{j+1,I}=x_{j,I}-A_{ii}(\cR)^{-1}\mathcal R_{\mathrm{skel}}(x_{j,I}),\qquad
  e_{j+1}\approx\big[I-A_{ii}(\cR)^{-1}A_{ii}(c_*)\big]e_j .
\end{equation}
This is the construction of \Cref{sec:iface} applied one level down: the
residual is evaluated exactly, so the fixed point is the same as that of
\eqref{eq:lchord}, and the frozen operator only sets the rate. The rate is
governed by the mismatch of the condensed operators rather than pointwise by
$\rho_i$; there is no a priori guarantee that it is smaller, but the
nonlinearity inside each leaf is resolved exactly, which is where the plain
chord iteration is expected to suffer at large $\Lambda$, when $c[\phi]$ varies
strongly within single leaves. A step costs one Newton solve per leaf and one
MUMPS solve, and no separate reconstruction: the leaf solutions are the field.
\end{remark}

\subsection{Choice of \texorpdfstring{$\cR$}{c\_ref}}
\begin{itemize}
  \item \emph{Constant} $\cR=\bar c$: simplest, but $c[\phi^*]$ varies like
    $e^{\phi^*/\tau}$ across $\Omega$, so at large $\Lambda$ no constant makes
    the local rates $\rho_i$ small.
  \item \emph{From a coarse solution} $\phi_c$: $\cR=c[\phi_c]$.
    Since $c[\phi]/c[\phi_c]=e^{(\phi-\phi_c)/\tau}$, to first order (iterates
    near $\phi^*$)
    \begin{equation}
      \rho_i\approx\sup_{\Omega_i}\frac{|\phi^*-\phi_c|}{\tau(\psi)},
    \end{equation}
    and the mismatch $S(\cR)-S(c_*)$ in \eqref{eq:irate} is small too. $\phi_c$
    is computed by the same algorithm on a cheap discretization (fewer slabs,
    lower order) with a constant $\bar c$, possibly with continuation in
    $\Lambda$.
\end{itemize}
The outer iterates are short vectors (one trace per interface), so Anderson
acceleration of \eqref{eq:ichord} costs almost nothing. So does an interface
Newton--Krylov method preconditioned by $P$ with finite-difference Jacobian
products: each product is one more evaluation of $\mathcal E$, which needs only
local solves.

\begin{algorithm}[h]
\caption{Fixed-point iteration on the interface system}\label{alg:chord}
\begin{algorithmic}[1]
\State Choose $\cR$; build the slab factorizations, HBS blocks and red--black preconditioner $P$ for $-\Delta+\cR$ \Comment{once}
\State $u_0\gets$ traces of a coarse solution, or the solution of the linear OMS system for $-\Delta+\cR$
\For{$k=0,1,2,\dots$}
  \For{each slab $i$} \Comment{independent, local}
    \State solve \eqref{eq:local} with data $u_{k,i\pm1}$ by \eqref{eq:lchord} (or \Cref{rem:leafelim})
    \State $\mathcal E_i(u_k)\gets x|_{\Gamma_i}$ \Comment{read from the skeleton solution}
  \EndFor
  \State $\mathcal R_k\gets u_k-\mathcal E(u_k)$; \textbf{if} $\|\mathcal R_k\|_\infty<\mathrm{tol}$ \textbf{then stop}
  \State $u_{k+1}\gets u_k-P\,\mathcal R_k$ \Comment{one red--black application; optionally Anderson}
\EndFor
\State the volumetric solution is given by the last local fields $\phi^{(i)}$ (no extra reconstruction)
\end{algorithmic}
\end{algorithm}

\begin{remark}
With exactly one local chord step per slab, each started from the previous
local field, and $P=(I+S_{\mathrm{ref}})^{-1}$ exact, \Cref{alg:chord}
gives $u_{k+1}=(I+S_{\mathrm{ref}})^{-1}b_{\mathrm{OMS}}(h_k)$, with $b_{\mathrm{OMS}}(h_k)$ the linear OMS
right-hand side for the body load $h_k=f^*-G[\phi_k]+\cR\phi_k$, $\phi_k$ the
current local fields. This is the
volumetric chord iteration $\phi\mapsto T(\phi)$ of \Cref{sec:nk}, except that
the local fields lag one interface update behind. Solving the local problems accurately removes the local
nonlinearity from the outer iteration: the outer rate then depends on the
interface mismatch $S(\cR)-S(c_*)$ instead of on $\rho$.
\end{remark}

%======================================================================
\section{Approach 2: Newton--Krylov preconditioned by \texorpdfstring{$M$}{M}}\label{sec:nk}
%======================================================================
\subsection{Newton system in preconditioned form}
This approach works with volumetric fields. Define the volumetric chord map
\begin{equation}\label{eq:T}
  T(\phi):=\SolveM\big(f^*-G[\phi]+\cR\,\phi,\ \phi^*\big),
\end{equation}
whose fixed point is the solution of \eqref{eq:mms}. Iterating
$\phi_{k+1}=T(\phi_k)$ converges with factor $\rho$, defined as in \eqref{eq:rho}
with the supremum over $\Omega$, by the same argument; Newton's method accelerates it.
Newton's method for \eqref{eq:mms} solves
\begin{equation}
  J_k\delta=-F(\phi_k),\qquad J_k=-\Delta+c_k,\quad c_k=c[\phi_k],\qquad \delta=0\ \text{on }\partial\Omega .
\end{equation}
Let $M_0^{-1}$ denote $\SolveM(\cdot,0)$. Two identities make this system
solvable with $M$ alone. For any $\phi$ with $\phi=\phi^*$ on $\partial\Omega$,
\begin{align}
  M_0^{-1}F(\phi)&=\phi-T(\phi), \label{eq:id1}\\
  M_0^{-1}J_k&=I+M_0^{-1}D_k,\qquad D_k=\diag(c_k-\cR). \label{eq:id2}
\end{align}
(For \eqref{eq:id1}: $M(\phi-T(\phi))=-\Delta\phi+\cR\phi-f^*+G[\phi]-\cR\phi=F(\phi)$,
and $\phi-T(\phi)$ vanishes on $\partial\Omega$.) The preconditioned Newton
system is therefore
\begin{equation}\label{eq:nk}
  \boxed{\;\big(I+M_0^{-1}D_k\big)\,\delta=T(\phi_k)-\phi_k,\qquad \phi_{k+1}=\phi_k+\delta .\;}
\end{equation}
The right-hand side costs one $\SolveM$. Each
operator application costs a pointwise multiplication by $c_k-\cR$ (the
diagonal update) and one $\SolveM$ with zero boundary data.

\subsection{Properties}
\begin{itemize}
  \item By \Cref{prop:maxp}, $\|M_0^{-1}D_k\|_\infty\le\sup_x|c_k-\cR|/\cR$:
    the preconditioned operator is a perturbation of the identity, and GMRES
    needs few iterations when $\cR$ is close to $c_k$, for instance when it comes from a coarse
    solution.
  \item The outer iteration converges quadratically near the solution; damping
    or a line search on $\|\phi-T(\phi)\|$ can be added for robustness, and the
    inner tolerance can follow an inexact-Newton (Eisenstat--Walker) rule.
  \item One Richardson step on \eqref{eq:nk} from $\delta=0$ gives
    $\delta=T(\phi_k)-\phi_k$, i.e.\ the volumetric chord step $\phi_{k+1}=T(\phi_k)$.
  \item $M_0^{-1}$ may be applied inexactly, e.g.\ with a single red--black
    application in place of the preconditioned interface GMRES. The outer
    Krylov method is then flexible GMRES.
  \item Krylov vectors are volumetric fields, stored on a non-overlapping set
    of leaf grids (each point owned by one slab).
\end{itemize}

\begin{algorithm}[h]
\caption{Newton--Krylov with a shared preconditioner}\label{alg:nk}
\begin{algorithmic}[1]
\State Build OMS, HBS and red--black for $M=-\Delta+\cR$ \Comment{once}
\State $\phi_0\gets\phi_c$ (or a chord iterate)
\For{$k=0,1,2,\dots$}
  \State $r_k\gets T(\phi_k)-\phi_k$ \Comment{one $\SolveM$}
  \State \textbf{if} $\|r_k\|_\infty<\mathrm{tol}$ \textbf{then stop}
  \State $c_k\gets c[\phi_k]$
  \State Solve $(I+M_0^{-1}D_k)\,\delta=r_k$ by (F)GMRES, with
         $v\mapsto v+\SolveM\big((c_k-\cR)\,v,\,0\big)$
  \State $\phi_{k+1}\gets\phi_k+\delta$
\EndFor
\end{algorithmic}
\end{algorithm}

%======================================================================
\section{What is reused and what is new}
%======================================================================
\begin{table}[h]
\centering
\begin{tabular}{@{}p{4.6cm}p{6.2cm}c@{}}
\toprule
Step & Uses & New factorization?\\
\midrule
Body-load reduction & leaf-level static condensation (small dense leaf solves; GPU) & no\\
Local solves & kept slab factorizations of $A_{ii}(\cR)$ & no\\
Leaf-level nonlinear elimination (\Cref{rem:leafelim}, optional) & small dense leaf Newton solves & leaf only\\
Interface solve & red--black application (\Cref{alg:chord}); GMRES on exact $I+S$ with red--black preconditioner (\Cref{alg:nk}) & no\\
Volumetric reconstruction & kept factorizations plus leaf solves & no\\
\bottomrule
\end{tabular}
\caption{Every iteration of both approaches consists of these steps and pointwise operations; the slab factorizations, HBS compressions and red--black preconditioner are never rebuilt.}
\end{table}

\paragraph{Requirements.} All slab factorizations are kept in memory
(\texttt{keepLU=True}), together with each slab's leaf operators
$L_{II,\ell}^{-1}$, $Y_\ell$, $D_\ell$ of \eqref{eq:leafops} (and $K_{II,\ell}$,
$K_{IF,\ell}$ for \Cref{rem:leafelim}), which are needed for body-load reduction
and reconstruction. For the five-patch
cylinder this means adding to the local slab solver (i) a body-load reducer
(per-patch leaf reduction, converted to the physical flux, glued and split into
interior and boundary parts) and (ii) a volumetric reconstruction with point
evaluation (patch classification, inverse patch map, interpolation on the leaf
grid). The OMS layer, the HBS compression and the red--black solver are used as
they are.

\begin{table}[h]
\centering
\begin{tabular}{@{}lll@{}}
\toprule
 & Interface fixed point (\Cref{alg:chord}) & Newton--Krylov (\Cref{alg:nk})\\
\midrule
State & interface traces $u$ (size $N_s n_c$) & volumetric fields\\
Global work per outer iteration & one red--black application & $1+n_{\mathrm{GMRES}}$ calls to $\SolveM$\\
Local work per outer iteration & local chord solves on every slab & leaf reduction and reconstruction per $\SolveM$\\
Convergence & linear; rate from $P$ and $S(\cR)-S(c_*)$ & quadratic (outer)\\
Sensitivity to $\cR$ & local rates $\rho_i$, interface mismatch & sets $n_{\mathrm{GMRES}}$\\
Extra memory & a few interface vectors & Krylov basis of volumetric fields\\
\bottomrule
\end{tabular}
\caption{Comparison of the two approaches.}
\end{table}

The interface iteration is the natural first choice: it keeps the OMS
unknowns as its state, needs no global volumetric storage, replaces the global
GMRES by one red--black application per outer iteration, and gets faster as
$\Lambda$ grows, because screening weakens the interface coupling. Both
approaches use exactly the same precomputed building blocks.

%======================================================================
\section{Verification}
%======================================================================
\begin{enumerate}
  \item \emph{Error against the manufactured solution} at points off the
    discretization grid, evaluated from the reconstructed field.
  \item \emph{Continuous residual.} For the reconstructed approximation
    $\tilde\phi$, evaluate $\eta=-\Delta\tilde\phi+G[\tilde\phi]-f^*$ inside
    leaves, the jumps of the normal derivative across leaf faces, and the
    boundary mismatch. Since $g$ is increasing in $s$, \Cref{prop:maxp} gives
    \begin{equation}
      \|\tilde\phi-\phi^*\|_\infty\le C_P\max_{\text{leaves}}|\eta|
      +C_\Gamma\max_{\text{faces}}\big|[\partial_n\tilde\phi]\big|
      +\max_{\partial\Omega}|\tilde\phi-\phi^*| ,
    \end{equation}
    with $C_P\le d^2/8$ for any slab of width $d$ containing $\Omega$ (for the
    torus above, $|x_3|\le\max(a_e,b_e)$ gives $d=2\max(a_e,b_e)$), and
    $C_\Gamma\approx C_P/h$ for leaf spacing $h$.
  \item \emph{Convergence studies:} the error and the residual bound against
    the polynomial order $p$ at fixed leaves, and outer and inner iteration
    counts against $\Lambda$, for constant and coarse-solution $\cR$.
\end{enumerate}

Unlike the interior Helmholtz problem, $M$ is coercive ($\cR\ge0$): there is no
resonance amplification of discretization errors, the interface operators have
low numerical rank, and the inner solves need few iterations.

\end{document}
